\Introduction

Официальная документация PostgreSQL служит рабочим источником для разработчиков, администраторов баз данных и инженеров эксплуатации. К ней обращаются при проектировании схем, проверке синтаксиса SQL-команд, настройке параметров сервера, диагностике репликации, анализе ограничений и сопровождении инсталляций в разных окружениях~\cite{postgres-docs,postgres-docs-16,postgres-docs-17,postgres-docs-18}. Сведения в документации публикуются по основным версиям PostgreSQL, а изменения поведения и набора возможностей фиксируются в политике версионирования и примечаниях к выпускам~\cite{postgres-versioning,postgres-release-root}. Поэтому инженерный ответ по документации должен быть не только связным, но и проверяемым по источникам выбранной версии.

Актуальность темы состоит в том, что на практике пользователь редко формулирует запрос как точное название раздела документации. Запрос обычно описывает задачу: проверить параметр, понять поведение команды, найти способ диагностики или получить пошаговое объяснение. Обычный поиск по сайту документации или по локальному корпусу HTML-документации возвращает набор страниц, но не формирует итоговый ответ и не контролирует, из какой версии взяты сведения. Использование больших языковых моделей без привязки к источникам также недостаточно надежно для инженерной эксплуатации, поскольку модель может смешать сведения разных версий или сформулировать ответ, который не подтверждается документацией~\cite{lewis2020rag,dpr2020,sbert2019,colbert2020,bert2018}.

Подход генерации с дополненным извлечением (Retrieval-Augmented Generation, RAG) позволяет разделить задачу на поиск релевантных фрагментов и генерацию ответа по найденному контексту~\cite{lewis2020rag}. В работе этот подход используется с ограничениями: корпус хранится локально, версия PostgreSQL передается как обязательный параметр, найденные источники возвращаются пользователю, а при недостатке подтверждений система не формирует уверенную рекомендацию. Такая постановка соответствует общим принципам проектирования информационных систем и требованиям к качеству программного обеспечения~\cite{sommerville2016,pressman2020,gost-r-iso-25010-2015}.

Объект исследования -- процесс предоставления пользователю справочной и обучающей информации по документации PostgreSQL с учётом выбранной версии.

Предмет исследования -- архитектурные, алгоритмические и интерфейсные решения веб-приложения, которое выполняет поиск фрагментов документации PostgreSQL, формирует ответ с источниками и сохраняет историю запросов.

Цель работы -- разработать веб-приложение, которое по запросу пользователя, выбранной версии PostgreSQL и режиму ответа извлекает релевантные фрагменты корпуса, проверяет достаточность подтверждений и формирует результат с источниками.

Для достижения цели поставлены следующие задачи:
\begin{compactlist}
  \item проанализировать предметную область, структуру документации PostgreSQL и проблему учёта версии;
  \item сравнить существующие подходы поиска по технической документации и определить ограничения ответов без источников;
  \item сформулировать функциональные и нефункциональные требования к веб-приложению;
  \item спроектировать архитектуру серверной части, модель данных, сценарий обработки запроса и алгоритмы работы с корпусом;
  \item реализовать хранение версий, документов, фрагментов, векторных представлений, истории запросов и источников в PostgreSQL;
  \item реализовать маршруты API, подготовку корпуса, поиск фрагментов, контроль версии, повторное ранжирование и проверку достаточности подтверждений;
  \item реализовать пользовательский интерфейс с вводом вопроса, выбором версии, выбором режима ответа, отображением результата, источников и истории;
  \item провести автоматическое и сценарное тестирование реализованных функций без заявления неподтверждённой численной оценки качества ранжирования.
\end{compactlist}

Работа состоит из четырех глав. В первой главе приведен анализ предметной области, документации PostgreSQL, существующих подходов поиска, рисков генеративных ответов, требований и постановки задачи. Во второй главе описаны проектные решения: архитектура веб-приложения, сценарий обработки запроса, диаграммы, модель данных и алгоритмы подготовки корпуса, поиска фрагментов, учёта версии, повторного ранжирования и проверки подтверждений. В третьей главе рассмотрена реализация серверной части, базы данных, маршрутов API, индексации, обработки запроса и пользовательского интерфейса. В четвертой главе приведены программа испытаний, автоматические тесты, сценарные проверки, результаты запуска тестов, сопоставление с требованиями и ограничения разработанного решения.

Методы исследования включают анализ предметной области, сравнительный анализ существующих подходов, структурно-функциональное моделирование, проектирование модели данных, проектирование программной архитектуры и функциональное тестирование. Для обоснования поисковой части учтены исследования в области информационного поиска, векторных представлений текста и вопросно-ответных систем~\cite{manning2008,bm25-prf,hnsw2016,faiss2017}. Для проектирования данных использованы классические работы по реляционной модели и базам данных~\cite{codd1970,ullman1988,date2004,garcia2008,silberschatz2020,elmasri2016}. Для архитектуры и сопровождения программной системы учтены подходы к разделению ответственности, эволюции кода и архитектурным стилям~\cite{martin2017,fowler2002,fowler2018,bass2021,larman2004,newman2021,richards2020}. Для проверки использованы общие принципы функционального тестирования~\cite{beizer1990,myers2011,kaner1999}.

Практическая значимость работы состоит в том, что разработанная система может использоваться как сервис поиска и объяснения материалов по документации PostgreSQL. Для разработчика приложений она помогает быстрее найти проверяемый ответ по SQL-командам и системным представлениям. Для администратора баз данных система полезна при проверке параметров, репликации, обслуживания и ограничений конкретной версии. Для инженера эксплуатации интерфейс с источниками и историей запросов упрощает повторную проверку рекомендаций при сопровождении PostgreSQL в локальной или контейнерной инфраструктуре~\cite{docker-docs,postgres-logical-replication-16,postgres-runtime-config-repl-16}.

В практической части реализована программная система на Python 3.12 с использованием FastAPI, SQLAlchemy, Alembic, PostgreSQL, pgvector, Pydantic, Beautiful Soup, httpx, Groq API и pytest~\cite{fastapi-docs,pydantic-docs,sqlalchemy-docs,alembic-docs,pgvector,beautifulsoup-docs,httpx-docs,groq-docs,pytest-docs}. Реализация поддерживает версии PostgreSQL 16, 17 и 18, краткий (\texttt{short}), подробный (\texttt{detailed}) и обучающий (\texttt{tutorial}) режимы ответа, локальную модель построения векторных представлений \texttt{BAAI/bge-m3} размерности 1024 и хранение истории пользовательских запросов.
